\uuid{Udel}
\exo7id{4605}
\titre{exo7 4605}
\auteur{quercia}
\organisation{exo7}
\datecreate{2010-03-14}
\isIndication{false}
\isCorrection{true}
\chapitre{Série entière}
\sousChapitre{Equations différentielles}
\cpge{mpsi-mp2i.primitives-equations-differentielles.b, pcsi.techniques-calcul-primitives.b, ptsi.techniques-calcul-primitives.b, mp-mpi.series-entieres.a, pc.series-entieres.a, psi.series-entieres.a, pt.series-entieres.a, mp-mpi.series-entieres.d, pc.series-entieres.c, psi.series-entieres.c, pt.series-entieres.c}
\module{Analyse}
\niveau{L2}
\difficulte{}

\contenu{
\question{
On note $\zeta_i(n) = \sum_{k=0}^\infty \frac1{(2k+1)^n}$
et $Z_i(x) = \sum_{n=1}^\infty \zeta_i(2n)x^n$.
En s'inspirant de l'exercice~\ref{recurzeta}
montrer que $Z_i$ vérifie l'équation différentielle~:
$2xZ_i'(x) - 2Z_i^2(x) - Z_i(x) = x\zeta_i(2)$.

Déterminer alors deux réels $\alpha$ et $\beta$ tels que
$T(x) = Z_i(x^2)/x$ soit égal à~$\alpha\tan\beta x$ sur~$]-1,1[$.
}
\reponse{
$\alpha=\pi/4$, $\beta=\pi/2$.
}
}
